\begin{python0}
from solutions import *; clear()
\end{python0}
\sectionthree{WARNING: SPOILERS!!!}

Here are the pseudocode (with explanations) for the last two ASCII art problems from the previous section.

Here's the first problem \ldots

\begin{ex}
Write a program that draws the following figure when the user enters 3
\begin{console}
***
* *
***
\end{console}
and when the user enters 4 it draws
\begin{console}
****
*  *
*  *
****
\end{console}
and here's the pseudocode with explanation:
%begin solution?

\tab[3em]{For n = 4:}\\
    \tab[5em]{4 stars}\\
    \tab[5em]{1 star, 2 spaces, 1 star}\\
    \tab[5em]{1 star, 2 spaces, 1 star}\\
    \tab[5em]{4 stars}\\    
\tab[3em]{For n = 5:}\\      
\tab[5em]{5 stars}\\
    \tab[5em]{1 star, 3 spaces, 1 star}\\
    \tab[5em]{1 star, 3 spaces, 1 star}\\
    \tab[5em]{1 star, 3 spaces, 1 star}\\
    \tab[5em]{5 stars}\\

In general, we have:\\
\tab[3em]{n stars}\\
\tab[3em]{n-2 of [1 star, n-2 stars, 1 star]}\\
\tab[3em]{n stars}\\

Therefore, the pseudocode is:

\tab[3em]{print n stars}\\
\tab[3em]{print newline}\\
\tab[3em]{for i = 1,..., n-2}\\
\tab[5em]{print 1 star}\\
\tab[5em]{print n-2 stars}\\
\tab[5em]{print 1 star}\\
\tab[5em]{print newline}\\
\tab[3em]{print n stars}\\

And after expanding some of the print statements we get:

\tab[3em]{//print n stars}\\
\tab[3em]{for i=1 ,..., n: print '*'}\\
\tab[3em]{print newline}\\
\tab[3em]{for i = 1, ..., n-2}\\
\tab[5em]{print '*'}\\
\tab[5em]{//print n-2 stars}\\
\tab[5em]{for j=1,..., n-2: print '*'}\\
\tab[5em]{print '*'}\\
\tab[5em]{print newline}\\
\tab[3em]{//print n stars}\\
\tab[3em]{for i=1, ..., n: print '*'}\\

\end{ex}
\begin{ex} Write a program that draws the following figure when
the user enters 3:
\begin{console}
  *
 ***
*****
\end{console}
and this when the user enters 4:
\begin{console}
   *
  ***
 *****
*******
\end{console}

\tab[3em]{When n = 3:}\\\\
     \tab[5em]{2 spaces, 1 star}\\
     \tab[5em]{1 space, 3 stars}\\
     \tab[5em]{0 space, 5 stars}\\\\
\tab[3em]{When n = 4:}\\\\
     \tab[5em]{3 spaces, 1 star}\\
     \tab[5em]{2 spaces, 3 stars}\\
     \tab[5em]{1 space, 5 stars}\\
     \tab[5em]{0 space, 7 stars}\\

There are two columns of numbers to handle. In other words, each line
(which is the body of a loop) has two numbers. The first column is easy - it's just from n-1 to 0. So the pseudocode looks like this:

    \tab[3em]{for i = n-1 to 0}\\
          \tab[5em]{i spaces, ??? stars}\\

What about the second? The relationship between the second number and i
is not so direct. But note that it starts with 1 and always goes up by 2
after each iteration. So the pseudocode looks like
     
      \tab[3em]{numStars = 1}\\
      \tab[3em]{for i = n-1 to 0}\\
              \tab[5em]{i spaces, numStars stars}\\
              \tab[5em]{numStars = numStars + 2}\\

In more details:

      \tab[3em]{numStars = 1}\\
      \tab[3em]{for i = n-1 to 0}\\
          \tab[5em]{print i spaces}\\
          \tab[5em]{print numStars stars}\\
          \tab[5em]{print newline}\\
          \tab[5em]{numStars = numStars + 2}\\

And after expanding some print statements we get

      \tab[3em]{numStars = 1}\\
      \tab[3em]{for i = n-1 to 0}\\
            \tab[5em]{//print i spaces}\\
            \tab[5em]{for j=1, \ldots, i: print ''}\\
            \tab[5em]{// print numStars stars}\\
            \tab[5em]{for j=1, \ldots, numStars: print '*'}\\
            \tab[5em]{print newline}\\
            \tab[5em]{numStars = numStars + 2}\\

\end{ex}
